\chapter*{Введение}
\addcontentsline{toc}{chapter}{Введение}
\label{ch:intro}

\textbf{JavaScript} --- один из трёх базовых языков веба (наряду с
\textbf{HTML} и \textbf{CSS}). HTML описывает структуру документа, CSS
задаёт его оформление, а JavaScript добавляет в страницу поведение:
реакцию на действия пользователя, асинхронный обмен данными с
серверами, изменение содержимого страницы без её перезагрузки. Без
JavaScript современный веб остановился бы на статических документах ---
любое интерактивное приложение (чат, веб-почта, интернет-магазин,
картографический сервис) построено вокруг JavaScript в браузере.

Предыдущие лабораторные работы (\textnumero\,10, \textnumero\,11,
\textnumero\,12) были посвящены проектированию интерфейса в графическом
редакторе \textbf{Figma}: студент создавал статичный макет, но не
оживлял его кодом. Текущая работа замыкает связку
<<\,макет \(\to\) код\,>>: на готовом HTML-каркасе с подключённым
\textbf{Bootstrap~5} реализуется одностраничное приложение
(<<\textit{Single Page Application}>>), которое получает данные с
публичного API через \texttt{fetch} и динамически собирает разметку
карточек и модальных окон.

В качестве темы выбран сайт-каталог супергероев вселенной
\textbf{Marvel}. Учебный API
\texttt{https://jsfree-les-3-api.onrender.com/characters} возвращает
JSON-массив из девятнадцати персонажей с полями \texttt{id},
\texttt{name} (русское имя), \texttt{nameor} (оригинальное имя),
\texttt{description}, \texttt{thumbnail} (URL обложки) и
\texttt{comics} (список комиксов с участием героя). Готовое
приложение публикуется на собственном домашнем сервере под обратным
прокси \textbf{Caddy} с автоматическим TLS-сертификатом Let's~Encrypt
и доступно по адресу
\url{https://marvel.server34.netcraze.club}.

Цель работы --- получить практические навыки разработки на JavaScript:
освоить базовые языковые конструкции, научиться работать с массивами,
объектами и асинхронными запросами, реализовать пользовательский
интерфейс на основе данных, полученных от внешнего API, и опубликовать
готовое приложение в интернете.

Задачи:
\begin{enumerate}
  \item Освоить базовые конструкции языка: переменные, типы данных,
        вывод в консоль (урок~1).
  \item Изучить работу с массивами, циклами и строковыми методами
        (\texttt{includes}, фильтрация) на примере мини-мессенджера
        (урок~2).
  \item Реализовать сайт-каталог героев Marvel: получить данные с
        учебного API через \texttt{fetch}, сформировать массив
        карточек и модальных окон, разместить их в DOM-дереве
        страницы (урок~3).
  \item Развернуть приложение на публичном хостинге, обеспечить
        работу по защищённому протоколу HTTPS (урок~4).
\end{enumerate}

\textbf{Стенд:}
\begin{itemize}
  \item \texttt{Node.js~18+} --- среда исполнения JavaScript вне
        браузера, используется для запуска консольных задач
        уроков~1 и~2;
  \item современный браузер (Chrome, Firefox, Safari) --- среда
        исполнения JavaScript на стороне клиента, используется для
        запуска сайта;
  \item \texttt{Bootstrap~5.0.1} (jsDelivr CDN) --- набор готовых
        CSS-компонентов и сетка, подключается одной строкой в
        \texttt{<head>};
  \item учебный REST-API
        \texttt{jsfree-les-3-api.onrender.com/characters} ---
        источник данных о персонажах Marvel;
  \item локальный HTTP-сервер \texttt{python3~-m~http.server}
        --- обход CORS-ограничений при открытии страницы по
        \texttt{file://};
  \item собственный домашний сервер (Debian~12, Docker Compose) с
        обратным прокси \texttt{Caddy} и автоматическим выпуском
        TLS-сертификатов Let's~Encrypt --- площадка для публикации
        готового сайта.
\end{itemize}

\endinput
